\documentclass{leveldoc}
\DocUses{foundations,deflection_plates}
\DocFingerprint{9647748f8bfc4567}

\begin{document}

\section{What the level asks}\label{sec:asks}

\leveldescription

One positive ion is launched along $+x$ from the left edge. The board starts
with no metal. The player places up to three plates, each at a potential from
the slider, inside the marked region. A positive plate pushes a positive ion
away. There are no charges and no magnets.

The ion must enter the detector. The shot has no acceptance window, no kinetic
window and no gate. A flight ends when the ion touches a plate, leaves the
world, or enters the detector (\xsecref{foundations}{sec:board}). A plate's
obstacle is its box plus a contact shell of thickness $\val{contact}$. The
ion's radius is $0$.

The reference places one positive plate above the launch line. The ion is
pushed downward and meets the detector's left side. With no plate the ion
flies straight, and the launch line passes above the detector. This document
uses the deflection plates' rectangle collocation. There is no second plate,
so there is no uniform gap: the check on the plate is its net charge seen from
the launch, and the empty board is an exact zero field. The flight is
integrated by the classical Runge--Kutta method.

\section{Setup}\label{sec:setup}

\setuptable

The speed of light is $c = \val{cLight}$. The ion has charge $q = \val{qPart}$,
mass $m = \val{mPart}$ and radius $0$, and it is launched with kinetic energy
$T_0 = \val{T0}$. The launch is at $(\val{xLaunch}, \val{yLaunch})$, straight
along $+x$. Then $\gamma_0 = \val{gamma0}$, $\abs{\vect{p}_0} = \val{p0}$ and
$v_0 = \val{v0} = \val{vOverC}\,c$ (\xeqref{foundations}{eq:launch}).

The level file does not list a plate size, so the plate is the engine's
default: length $\val{plateL}$ along $x$, thickness $\val{plateT}$ and height
$\val{plateH}$. The reference plate is centred at
$(\val{xPlateC}, \val{yPlateC})$, at potential $V = \val{vPlate}$, and it
occupies $\val{xPlate0} \le x \le \val{xPlate1}$,
$\val{yPlate0} \le y \le \val{yPlate1}$. The slider runs from $-\val{vSlider}$
to $\val{vSlider}$.

The detector is $\val{xDetLo} \le x \le \val{xDetHi}$,
$\val{yDetLo} \le y \le \val{yDetHi}$. The world is $-1 \le x \le \val{xWorld}$,
$-1 \le y \le \val{yWorld}$. The reported time step is $h = \val{H}$. The mesh
uses rectangles of edge $\val{panelFine}$, $\val{nPanels}$ of them on the
upper half of the one plate.

\begin{figure}[ht]
\centering
\begin{tikzpicture}
\begin{axis}[axis equal image, width=0.92\linewidth,
  xmin=-1.5, xmax=32.5, ymin=-2, ymax=22.5,
  xlabel={$x$}, ylabel={$y$}, tick label style={font=\small}]
  \draw[black!40, dashed] (-1,-1) rectangle (\rawval{xWorld},\rawval{yWorld});
  \draw[black!25] (0,0) rectangle (30,20);
  \fill[green!50!black, fill opacity=0.18]
    (\rawval{xDetLo},\rawval{yDetLo}) rectangle (\rawval{xDetHi},\rawval{yDetHi});
  \draw[green!40!black]
    (\rawval{xDetLo},\rawval{yDetLo}) rectangle (\rawval{xDetHi},\rawval{yDetHi});
  \fill[black!20]
    (\rawval{xPlate0},\rawval{yPlate0}) rectangle (\rawval{xPlate1},\rawval{yPlate1});
  \draw[black!70]
    (\rawval{xPlate0},\rawval{yPlate0}) rectangle (\rawval{xPlate1},\rawval{yPlate1});
  \addplot[orange!80!black] table {pathbare.dat};
  \addplot[very thick, blue!60!black] table {pathref.dat};
  \draw[-{Stealth[length=5pt]}, thick] (\rawval{xLaunch},\rawval{yLaunch}) -- (\rawval{arrowX},\rawval{arrowY});
  \fill (\rawval{xLaunch},\rawval{yLaunch}) circle[radius=0.08];
\end{axis}
\end{tikzpicture}
\caption{The plate and the two flights. The gray rectangle is the reference
plate, above the launch line; the green box is the detector. Blue is the
reference flight: the positive plate pushes the ion down, and the ion meets
the left side of the detector. Orange is the same launch with no plate: the
ion stays on the launch line and leaves through the right edge of the world,
above the detector. The blue path stays outside the plate.}
\label{fig:board}
\end{figure}

\section{Principles and model}\label{sec:model}

\begin{principles}
  \item The relativistic equation of motion \xeqref{foundations}{eq:eom}, the
    launch \xeqref{foundations}{eq:launch}, and energy in a static field
    \xeqref{foundations}{eq:energy}. There is no magnetic field.
  \item A player plate is a metal box at the potential the player sets
    (\physref{2.7}). Its surface charge is the deflection plates' collocation
    \xeqref{deflection_plates}{eq:collocation}, with the rectangle integral
    \xeqref{deflection_plates}{eq:rectangle}.
  \item With no metal and no charges the field is zero, so the motion is
    uniform. That flight is exact.
  \item The board's rules (\xsecref{foundations}{sec:board}) and the classical
    Runge--Kutta method (\xsecref{foundations}{sec:numerics}).
  \item The radiated power \xeqref{foundations}{eq:larmor}, summed by Simpson's
    rule \xeqref{foundations}{eq:simpson}. It is not put back into the force.
\end{principles}

The reference has one plate at $V = \val{vPlate}$ and no other charge. The
state obeys \xeqref{foundations}{eq:eom} with $\vect{B} = \vect{0}$ until the
first ending event.

\begin{assumption}{A static plate}\label{as:static}
The potential equals the slider value at once. The launch speed is
$\val{vOverC}\,c$. The game uses the same static plate (\physref{2.7}).
\end{assumption}

\begin{assumption}{No radiation reaction}\label{as:radiation}
The radiated energy is not taken from the flight. On the reference flight it
is a fraction $\val{RefRadFrac}$ of $T_0$.
\end{assumption}

\begin{assumption}{No image of the ion}\label{as:image}
The ion's own surface charge is left out, as on the deflection plates and as
in the game (\physref{2.7}). The bound $q^2/d^2$, with $d$ floored at
$10^{-3}$ and taken relative to $\max(\abs{q\vect{E}}, T_0)$, reaches
$\val{RefImage}$ on the reference samples. The empty board has no metal. The
engine asks for the ratio to stay below $10^{-10}$.
\end{assumption}

\begin{assumption}{Equal rectangles}\label{as:panels}
The mesh is equal rectangles of edge $\val{panelFine}$. Replacing the edge by
$\val{panelCoarse}$ moves an end state by $\val{errMesh}$. The empty board has
no panels, so that flight does not depend on the mesh.
\end{assumption}

\begin{assumption}{The fixed step}\label{as:step}
The reported step is $h = \val{H}$. Halving it estimates an end state by
$\val{errStep}$. Ten times the panel-size estimate $\val{errMesh}$ lies above
$10^{-9}$, and it is larger than the step-halving estimate $\val{errStep}$, so
the tolerance on a state is that product, $\val{tolFlight}$. Between samples
the ion moves at most $\val{stepReach}$. The sampled clearance of the plate is
$\val{RefDmin}$.
\end{assumption}

\section{Derivation}\label{sec:derivation}

\subsection{One plate}\label{sec:plate}

A player plate is an electrode. The rectangles, the mirror and the collocation
are \xeqref{deflection_plates}{eq:rectangle} and
\xeqref{deflection_plates}{eq:collocation}, with the plate potential on the
right-hand side and with no external charge. The reported mesh has
$\val{nPanels}$ upper rectangles. The residual at the centroids is
$\val{solveRes}$. On the face toward the beam, off the centroids, the
potential differs from $V$ by at most $\val{surfRes}$.

There is no second plate, so the uniform-gap field does not apply. Two checks
replace it. On the launch line under the plate's centre the panel field is
$\val{eUnder}$, downward, and the potential is $\val{phiUnder}$. The net charge
of the plate and its mirror is $\val{qNet}$. At the launch, a distance
$\val{rLaunch}$ from the plate's centre, that charge would produce the Coulomb
potential $\val{phiMono}$. The panel potential there is $\val{phiLaunch}$, and
the relative difference is $\val{monoRel}$. That is the size of treating the
plate as a point at the launch. The march uses the panels, because the ion
passes $\val{RefDmin}$ from the metal, where the point charge is not the field.

With the plate omitted there is no charge and no metal, so $\vect{E} = \vect{0}$.
The field sampled on the launch line is $\val{eBare}$.

\subsection{The flight}\label{sec:flight}

The diagnostic is \xeqref{deflection_plates}{eq:flight-energy}. The state is
advanced by \xeqref{foundations}{eq:rk4} with step $h = \val{H}$. Events are
located as on the deflection plates: thirty-two Hermite samples, then sixty
bisections, obstacles before the world, and the world before the detector. The
step-halving estimate divides the change by $15$; the panel-size estimate is
the change itself when the edge doubles.

\section{The reference solution by hand}\label{sec:hand}

\begin{procedure}
  \item\label{st:launch} \emph{The launch} \xeqref{foundations}{eq:launch}.
    $\gamma_0 = \val{gamma0}$, $\abs{\vect{p}_0} = \val{p0}$, so
    $\vect{p}_0 = (\val{p0}, 0)$ and $v_0 = \val{v0}$.
  \item\label{st:plate} \emph{The plate.} Centre
    $(\val{xPlateC}, \val{yPlateC})$, potential $\val{vPlate}$, size
    $\val{plateL}$ by $\val{plateT}$ by $\val{plateH}$. Solve
    \xeqref{deflection_plates}{eq:collocation}. The centroid residual is
    $\val{solveRes}$ and the off-centroid discrepancy is $\val{surfRes}$. The
    net charge is $\val{qNet}$. At the launch the panel potential and the
    point-charge potential differ by the relative amount $\val{monoRel}$.
  \item\label{st:march} \emph{The march.} Advance
    \xeqref{deflection_plates}{eq:flight-energy} and
    \xeqref{foundations}{eq:eom} by \xeqref{foundations}{eq:rk4} with step
    $\val{H}$. The state every five time units:
    \begin{center}\small\val{states}\end{center}
  \item\label{st:end} \emph{The arrival.} The ion meets the \val{RefFace} face
    of the detector at $(\val{RefX}, \val{RefY})$, at $t = \val{RefT}$, with
    momentum $(\val{RefPx}, \val{RefPy})$, at an angle $\val{RefAng}$ degrees
    from $+x$. The abscissa is the detector's left edge, and the ordinate lies
    between $\val{yDetLo}$ and $\val{yDetHi}$. The run takes $\val{RefSteps}$
    steps. The path only falls: its highest point is the launch, and its
    lowest is $\val{RefYmin}$.
  \item\label{st:check} \emph{The checks.} $W$ drifts by at most
    $\val{RefWdev}$. The step-halving estimate is $\val{RefErr}$, and the
    larger of the two flights is $\val{errStep}$. The panel-size estimate is
    $\val{errMesh}$. The sampled clearance is $\val{RefDmin}$, outside the
    shell $\val{contact}$. The image-force ratio reaches $\val{RefImage}$.
  \item\label{st:rad} \emph{The radiation.} Simpson's rule on steps of $h/2$
    gives $\val{RefRad}$, a fraction $\val{RefRadFrac}$ of $T_0$. Halving that
    step again changes the integral by the relative amount $\val{RefRadErr}$.
    This integral is not a row of the comparison.
\end{procedure}

\begin{result}[The reference flight]
The ion enters the detector through its left side at
$(\val{RefX}, \val{RefY})$ at $t = \val{RefT}$, with momentum
$(\val{RefPx}, \val{RefPy})$. The level is solved.
\end{result}

\section{Comparison with the game}\label{sec:comparison}

\comparisontable

Eight comparisons are in the table: the reference outcome, time and end
position, and the outcome, time and end position with no plate. The tolerance
on a state is $\val{tolFlight}$. Ten times the panel-size estimate
$\val{errMesh}$ lies above $10^{-9}$, and it is larger than the step-halving
estimate $\val{errStep}$, so the tolerance is that product. The tolerance on a
state was fixed from these estimates before the comparison.

The largest absolute difference is the reference ordinate,
$\val{cmpRefYDiff}$. It lies above the panel-size estimate $\val{errMesh}$ and
under the tolerance. The reference time differs by $\val{cmpRefTDiff}$, and
the bare time by $\val{cmpBareTDiff}$. The outcomes agree.

The reference abscissa, the bare abscissa and the bare ordinate agree exactly.
The reference abscissa is the detector's left edge, the bare abscissa is the
world's right edge, and the bare ordinate is the launch line.

\section{Other solutions}\label{sec:others}

No plate, which is the empty placement in \code{placements.json}, does not
solve the level. The field is zero, and the ion stays on the launch line. It
leaves through the \val{BareFace} edge of the world at
$(\val{BareX}, \val{BareY})$, at $t = \val{BareT}$, with momentum
$(\val{BarePx}, \val{BarePy})$. The ordinate $\val{BareY}$ is above the
detector, which ends at $\val{yDetHi}$. The energy drifts by $\val{BareWdev}$,
and the run takes $\val{BareSteps}$ steps.

The reference works because the plate is positive and sits above the launch, so
the ion is pushed down from $\val{yLaunch}$ to $\val{RefY}$, inside the
detector. The slider's half-range is $\val{vSlider}$. This document integrates
the reference plate and the empty board, not the other slider values and not
a second or third plate.

\section{What the game computes}\label{sec:engine}

A player plate is the same boundary-element electrode as a level plate, with
the default size when the level file omits it (\physref{2.7}). This document
collocates \xeqref{deflection_plates}{eq:collocation} on equal rectangles.
The end states move by $\val{errMesh}$ when that edge doubles, and the
reference ordinate in \S\ref{sec:comparison} lies above that estimate and under
ten times it. The empty board is exact on both sides: there is nothing to
mesh.

The ion's image force is omitted on both sides. The bound sampled on the
reference flight reaches $\val{RefImage}$, under the engine's limit
$10^{-10}$. The flight in the engine is the adaptive method of \physref{5} and
\physref{6}, verified as in \physref{7}. Our step estimate is $\val{errStep}$.
Radiation is integrated and not applied to the momentum
(\xsecref{foundations}{sec:radiation}, \physref{8}).

\section{Further reading}\label{sec:reading}

The rectangle integral and the collocation are derived for the deflection
plates (\xeqref{deflection_plates}{eq:rectangle},
\xeqref{deflection_plates}{eq:collocation}). The motion and the Runge--Kutta
step are in the Foundations (\xsecref{foundations}{sec:numerics}). The
engine's player plates are \physref{2.7}.

\levelfileappendix

\end{document}
